% TEMPLATE-MILC-v3.tex - plantilla maestra UNICA de las guias MILC v3.
%
% La usan las cuatro familias de guias, cada una con su driver:
%   · Grados 6-11          -> scripts/build-guias-g11.py
%   · Bebras               -> scripts/build-guias-bebras.py
%   · Semillero            -> scripts/build-guias-semillero.py
%   · Territorio interior  -> scripts/build-guias-territorio-interior.py
%
% Antes habia tres copias casi identicas (~400 lineas iguales, 6 distintas):
% cada mejora habia que aplicarla tres veces, y en la practica no se hacia
% (el motor de recursos vivia solo en dos de ellas). Lo que varia entre
% familias esta parametrizado en 6 placeholders que cada driver llena
% (se nombran aqui SIN su sintaxis real: el builder reemplaza por texto plano,
% tambien dentro de los comentarios, y un valor multilinea rompe el preambulo):
%
%   HEADER_IZQ        encabezado izquierdo de pagina
%   HEADER_DER        encabezado derecho de pagina
%   PORTADA_ETIQUETA  rotulo sobre el numero grande (GRADO/MOMENTO/MODULO)
%   PORTADA_SERIE     linea de serie de la portada
%   PORTADA_ID        identificador de la guia en la portada
%   PORTADA_CIRCULO   el circulito de grado (°), o vacio si no aplica
%   PDF_TITULO        titulo en texto plano (sin \textbf ni comillas TeX) para
%                     los metadatos del PDF (/Title); lo produce lib_xelatex.texto_pdf
%
% Composicion (auditoria 2026-09-03): las bandas de estacion piden espacio
% (\Needspace*) y prohiben el salto justo despues (\nopagebreak), asi no quedan
% huerfanas al pie; las cajas largas (softbox, drawbox, infoband, puentebox) son
% `breakable`, asi una caja de media pagina ya no arrastra todo a la siguiente
% dejando la anterior al 40 %. Todo texto sobre mostaza o verde va en milcNegro
% (blanco sobre #E5B400 da 1,93:1 y sobre #58A923 2,95:1; ninguno pasa WCAG AA).
% El resto de claves las documenta content/guias/_SCHEMA.md. El script valida
% que no queden placeholders sin reemplazar antes de escribir el archivo final.

% Etiquetado PDF (latex-lab): /Lang, estructura y texto alternativo de las
% figuras, para lectores de pantalla. Debe ir ANTES de \documentclass.
\DocumentMetadata{lang=es-CO,tagging=on}
\documentclass[11pt,letterpaper]{article}
% headheight=24pt: la cabecera derecha lleva dos lineas (identidad de la guia
% y estacion actual). headsep se acorta en la misma medida para que el bloque
% de texto quede exactamente donde estaba con headheight=12pt.
\usepackage[margin=2.0cm,headheight=24pt,headsep=13pt]{geometry}
\usepackage[table]{xcolor}
\usepackage{fontspec}
\usepackage[spanish]{babel}
\usepackage{tikz}
% Librerías TikZ para los diagramas de las guías (bloque `recursos:`):
%  · babel  → babel en español vuelve activos caracteres como «>», y TikZ los usa
%             en sus opciones (>=stealth, ->). Sin esta librería, cualquier
%             diagrama con flechas revienta con «\language@active@arg> has an
%             extra }». Debe ir ANTES de que se dibuje nada.
%  · positioning        → sintaxis «below left=2pt and 3pt».
%  · decorations.pathreplacing → llaves ({brace}) para señalar rangos.
\usetikzlibrary{babel,positioning,decorations.pathreplacing}
\usepackage{graphicx}
% Circuitos: el trabajo de electrónica se hace hoy con micro:bit y sus sensores
% integrados (MakeCode, sin cableado), así que no se necesitan esquemáticos.
% Si algún día entran montajes con protoboard, basta descomentar esta línea:
% los diagramas .tex/.tikz declarados en `recursos:` ya se incrustan solos.
% \usepackage{circuitikz}
\usepackage[most]{tcolorbox}
\usepackage{tabularx}
\usepackage{array}
\usepackage{fancyhdr}
\usepackage{titlesec}
\usepackage[document]{ragged2e}  % bandera a la derecha en todo el cuerpo
\usepackage{enumitem}
\usepackage{microtype}
\usepackage{etoolbox}
\usepackage{needspace}
% hyperref al final del preambulo: metadatos del PDF (titulo, autor, idioma,
% familia), marcadores por estacion (\pdfbookmark) y enlaces sin recuadro.
\usepackage[hidelinks,
  pdftitle={Depuración — una hipótesis, una prueba},
  pdfauthor={PhD. Álvaro Cárdenas Orozco},
  pdfsubject={Tecnología e Informática · Grado 8 · Guía 17 · MILC v3},
  pdflang={es-CO},
  bookmarksopen=true,bookmarksnumbered=false]{hyperref}

% Helvetica Neue no trae la flecha U+2192, asi que cada «→» escrito literalmente
% en el YAML se compilaba a NADA, en silencio (462 flechas en 61 guias: rutas del
% tipo «paleta → tipografia → lockup» salian rotas). Mapeamos el caracter a su
% version matematica, que si tiene glifo. Asi el autor puede seguir escribiendo
% «→» comodamente en el YAML.
\usepackage{newunicodechar}
\newunicodechar{→}{\ensuremath{\rightarrow}}
% Mismo problema con los simbolos de marca y vinetas: el «✓» de «una palomita ✓»
% salia como cuadrito vacio en el PDF de todas las guias que lo usaban.
\usepackage{pifont}
\newunicodechar{✓}{\ding{51}}
\newunicodechar{✗}{\ding{55}}
\newunicodechar{✔}{\ding{52}}
\newunicodechar{•}{\textbullet}
\newunicodechar{≥}{\ensuremath{\geq}}
\newunicodechar{≤}{\ensuremath{\leq}}

\setmainfont{Helvetica Neue}
\setsansfont{Helvetica Neue}
\setlength{\parindent}{0pt}
\setlength{\parskip}{6pt}
% Sin particion de palabras: en 6.º-7.º una palabra cortada por guion al final
% de linea («Comuni-cacion») frena la lectura mucho mas que una linea corta.
\hyphenpenalty=10000
\exhyphenpenalty=10000
\pretolerance=2000
\tolerance=3000
\setlength{\emergencystretch}{3em}
\linespread{1.10}
\renewcommand{\arraystretch}{1.25}
\setlist{leftmargin=5mm,itemsep=1mm,topsep=1mm}
\newcolumntype{Y}{>{\RaggedRight\arraybackslash}X}

\definecolor{milcMagenta}{HTML}{E6007E}
\definecolor{milcVino}{HTML}{5A0038}
\definecolor{milcTurquesa}{HTML}{0093A5}
\definecolor{milcVerde}{HTML}{58A923}
\definecolor{milcMostaza}{HTML}{E5B400}
\definecolor{milcOcre}{HTML}{B86600}
\definecolor{milcNegro}{HTML}{241816}
\definecolor{milcGris}{HTML}{F7F7F4}
\definecolor{milcLinea}{HTML}{E8E2DD}
\definecolor{milcGradePrimary}{HTML}{FF6600}
\definecolor{milcGradeDark}{HTML}{9A3E00}
\definecolor{milcGradeSoft}{HTML}{FFE4D0}
\definecolor{milcGradeAccent}{HTML}{CFE7FF}
\definecolor{milcGradeText}{HTML}{FFFFFF}
\color{milcNegro}

\pagestyle{fancy}
\fancyhf{}
% Cabecera de dos lineas: identidad de la guia arriba y, a la derecha y en
% gris, la estacion en curso (\markright desde \milcestacion). Antes de la
% primera estacion la segunda linea va vacia; el \strut mantiene la altura
% para que la linea de arriba no baile entre paginas.
\lhead{\small Tecnología e Informática · Grado 8\\[-1pt]{\footnotesize\strut}}
\rhead{\small Guía 17 · MILC v3\\[-1pt]{\footnotesize\strut\color{milcNegro!65}\rightmark}}
\cfoot{\small\thepage}
\renewcommand{\headrulewidth}{0pt}

% Sobre mostaza (#E5B400) y verde (#58A923) el texto va en milcNegro; sobre el
% resto de la paleta, en blanco. Se usa dentro de fonttitle/fontupper, donde un
% \color posterior manda sobre coltitle.
\newcommand{\milctintasobre}[1]{%
  \ifboolexpr{test{\ifstrequal{#1}{milcMostaza}} or test{\ifstrequal{#1}{milcVerde}}}%
    {\color{milcNegro}}{\color{white}}}

\newtcolorbox{titlebox}[1]{
  enhanced,colback=#1,colframe=#1,arc=7mm,boxrule=0pt,
  left=7mm,right=7mm,top=3mm,bottom=3mm,
  before skip=8pt,after skip=8pt,
  fontupper=\sffamily\bfseries\Large\color{white}
}

\newtcolorbox{softbox}[3]{
  enhanced,breakable,pad at break=3mm,
  colback=#2,colframe=#1,borderline west={4pt}{0pt}{#1},
  arc=2mm,boxrule=.4pt,left=4mm,right=4mm,top=3mm,bottom=3mm,
  before skip=6pt,after skip=8pt,title={#3},coltitle=white,
  colbacktitle=#1,fonttitle=\sffamily\bfseries\milctintasobre{#1},fontupper=\RaggedRight,
  attach boxed title to top left={xshift=3mm,yshift=-2mm},
  boxed title style={arc=2mm, boxrule=0pt}
}

\newtcolorbox{drawbox}[2]{
  enhanced,breakable,pad at break=4mm,
  colback=white,colframe=#1,arc=4mm,boxrule=1pt,
  left=5mm,right=5mm,top=4mm,bottom=4mm,before skip=8pt,after skip=8pt,
  title={#2},coltitle=white,colbacktitle=#1,fonttitle=\sffamily\bfseries\milctintasobre{#1},
  fontupper=\RaggedRight,
  attach boxed title to top left={xshift=4mm,yshift=-2mm},
  boxed title style={arc=3mm, boxrule=0pt}
}

\newtcolorbox{infoband}[1]{
  enhanced,breakable,pad at break=2mm,colback=#1!8,colframe=#1!50,arc=2mm,boxrule=.5pt,
  left=4mm,right=4mm,top=2mm,bottom=2mm,before skip=4pt,after skip=4pt,
  fontupper=\small\RaggedRight
}

% Puente narrativo entre secciones (contrato editorial Conéctate).
% Da continuidad: "Acabas de X. Ahora vas a Y porque Z."
% Visualmente: banda izquierda en color del grado, texto en itálica pequeña.
\newtcolorbox{puentebox}{
  enhanced,breakable,pad at break=2mm,colback=milcGradeSoft,colframe=milcGradeDark,
  borderline west={3pt}{0pt}{milcGradeDark},
  arc=0mm,boxrule=0pt,left=5mm,right=4mm,top=2.5mm,bottom=2.5mm,
  before skip=4pt,after skip=8pt,
  fontupper=\itshape\small\color{milcNegro}\RaggedRight
}
% La flecha va en modo matematico a proposito: Helvetica Neue NO tiene el glifo
% U+2192, asi que \textrightarrow se compilaba a NADA («Missing character: There
% is no → in font Helvetica Neue») y el puente salia sin flecha. En modo
% matematico la flecha viene de la fuente matematica y si se dibuja.
\newcommand{\puente}[1]{%
  \begin{puentebox}%
    \textcolor{milcGradeDark}{$\rightarrow$}\hspace{2mm}#1%
  \end{puentebox}%
}

\newcommand{\linead}[1]{\noindent\color{milcLinea}\rule{#1}{0.5pt}\color{milcNegro}}
\newcommand{\respuesta}[1]{\par\vspace{2mm}\foreach \n in {1,...,#1}{\noindent\color{milcLinea}\rule{\linewidth}{0.5pt}\par\vspace{5mm}}\color{milcNegro}}
\newcommand{\checkbox}{\textcolor{milcGradeDark}{\fbox{\phantom{x}}}\hspace{5pt}}

% ─── Listas aireadas para las guías socioemocionales ────────────────────
% Componer las verificaciones como «\checkbox texto\\» deja el texto que salta
% de línea debajo del cuadrito y apelmaza el bloque. Estos entornos alinean la
% continuación y airean los ítems. Los usa Territorio Interior; cualquier
% familia puede hacerlo.
\newlist{checklista}{itemize}{1}
\setlist[checklista]{
  label=\textcolor{milcGradeDark}{\fbox{\phantom{x}}},
  leftmargin=9mm, labelsep=3.2mm, itemsep=2.4mm,
  topsep=2.5mm, parsep=0pt, partopsep=0pt, before=\RaggedRight
}
\newlist{pasoslista}{enumerate}{1}
\setlist[pasoslista]{
  label=\textcolor{milcGradeDark}{\sffamily\bfseries\arabic*.},
  leftmargin=8mm, labelsep=3mm, itemsep=2.8mm,
  topsep=1mm, parsep=0pt, partopsep=0pt, before=\RaggedRight
}

\newcommand{\phasecard}[4]{
\begin{tcolorbox}[enhanced,colback=#3,colframe=#2,arc=3mm,boxrule=.5pt,height=3.05cm,valign=center,halign=center,left=1.5mm,right=1.5mm,top=2mm,bottom=2mm]
{\sffamily\bfseries\fontsize{22}{24}\selectfont\textcolor{#2}{#1}}\par
{\sffamily\bfseries\small #4}
\end{tcolorbox}
}

% ── Piezas del rediseno 2026 (legibilidad 6.º-11.º) ───────────────────
% Banda de estacion: reemplaza el titlebox macizo. Titulo a la izquierda,
% dato practico (tiempo/modalidad) a la derecha, en una sola linea.
% Nunca huerfana: pide 9 lineas libres (o salta de pagina rellenando con
% \vfill) y prohibe el salto justo despues de la banda. Nueve y no seis:
% la caja partible que sigue necesita ~80pt (titulo adosado, relleno y dos
% lineas) para arrancar en la misma pagina; con menos, tcolorbox la manda
% entera a la siguiente y la banda se queda sola. El penalty 10000 va
% en `after`, ANTES del espacio posterior: un `after skip` (aunque sea 0pt)
% es un \vskip tras la caja, y ese glue es un punto de corte legal que un
% \nopagebreak escrito despues de \end{tcolorbox} ya no cierra. Cada banda
% deja marcador PDF y actualiza la cabecera derecha.
\newcounter{milcestacion}
\newcommand{\milcestacion}[2]{%
\Needspace*{9\baselineskip}%
\stepcounter{milcestacion}%
\pdfbookmark[1]{#2}{milcestacion\themilcestacion}%
\markright{#2}%
\begin{tcolorbox}[enhanced,colback=#1,colframe=#1,arc=2mm,boxrule=0pt,
  left=5mm,right=5mm,top=2.6mm,bottom=2.6mm,before skip=11pt,
  after={\par\nopagebreak\vskip7pt}]
{\sffamily\bfseries\fontsize{15}{18}\selectfont\milctintasobre{#1}#2}
\end{tcolorbox}}

% Tarjeta de paso numerada: sustituye las tablas «Pilar / Aplicacion», que
% eran formato de consulta docente, no de estudiante. El numero va en un
% circulo sobre el borde para que la secuencia se lea de un vistazo.
\newcommand{\milcpaso}[3]{%
\Needspace{4\baselineskip}%
\begin{tcolorbox}[enhanced,colback=white,colframe=milcLinea,arc=1.5mm,boxrule=.9pt,
  left=14mm,right=4mm,top=3mm,bottom=3mm,before skip=4pt,after skip=4pt,
  fontupper=\RaggedRight,
  overlay={\node[anchor=north west,circle,fill=#2,text=white,inner sep=0pt,
    minimum size=7.4mm,font=\sffamily\bfseries\small]
    at ([xshift=3.6mm,yshift=-3.2mm]frame.north west) {#1};}]
#3
\end{tcolorbox}}

% Casilla marcable de tamano util con lapiz (la anterior era \fbox{\phantom{x}},
% de alto variable segun el texto que la seguia).
\newcommand{\milcmarca}{\framebox[3.8mm]{\rule{0pt}{3.4mm}}}

% Fila marcable de lista suelta (checks de la estacion 1).
\newcommand{\milccheckrow}[1]{%
\par\vspace{1.8mm}\noindent
\makebox[6.4mm][l]{\raisebox{-.55mm}{\textcolor{milcNegro}{\milcmarca}}}%
\parbox[t]{\dimexpr\linewidth-6.4mm\relax}{\RaggedRight #1}\par}

% Celda marcable dentro de la rubrica (tabla de autoevaluacion). Misma
% sangria francesa, pero pensada para vivir dentro de una columna X.
\newcommand{\milcopcion}[1]{%
\makebox[5.2mm][l]{\raisebox{-.55mm}{\textcolor{milcNegro}{\milcmarca}}}%
\parbox[t]{\dimexpr\linewidth-5.2mm\relax}{\RaggedRight #1}}

% ── Recursos visuales de la guía ──────────────────────────────────────
% Declarados en el bloque `recursos:` del YAML y emitidos por el builder.
% \guiaFigura   → imagen rasterizada o PDF (foto, carta celeste, montaje).
% \guiaDiagrama → fuente TikZ/circuitikz (.tex) incrustada como vector:
%                 es la vía para circuitos y esquemáticos editables.
% [#1] = texto alternativo (alt) · #2 = ruta relativa al .tex · #3 = pie
% El alt llega al PDF etiquetado (/Alt) y lo lee el lector de pantalla.
\newcommand{\guiaFigura}[3][]{%
\begin{center}
\includegraphics[width=0.88\linewidth,height=0.38\textheight,keepaspectratio,alt={#1}]{#2}\\[1.5mm]
{\footnotesize\itshape\textcolor{milcNegro!75}{#3}}
\end{center}
}

\newcommand{\guiaDiagrama}[2]{%
\begin{center}
\input{#1}\\[1.5mm]
{\footnotesize\itshape\textcolor{milcNegro!75}{#2}}
\end{center}
}

\begin{document}
\thispagestyle{empty}

% ── Portada ───────────────────────────────────────────────────────────
% Antes: un rectangulo de color a pagina completa. Se veia bien en pantalla,
% pero en la fotocopiadora del colegio se comia el toner y salia gris sucio.
% Ahora la banda ocupa el tercio superior y el resto va en blanco.
\begin{tikzpicture}[remember picture,overlay]
  \path[fill=milcGradePrimary]
    (current page.north west) rectangle ([yshift=-10.6cm]current page.north east);

  \node[anchor=north west,text=milcGradeText] at ([xshift=2.0cm,yshift=-1.55cm]current page.north west)
    {\sffamily\bfseries\fontsize{12}{14}\selectfont GRADO};
  \node[anchor=north west,text=milcGradeText,opacity=.85] at ([xshift=2.0cm,yshift=-2.35cm]current page.north west)
    {\sffamily\bfseries\fontsize{8.5}{10}\selectfont SERIE GUÍAS · TECNOLOGÍA E INFORMÁTICA};
  \node[anchor=north east,text=milcGradeText,opacity=.85,align=right] at ([xshift=-2.0cm,yshift=-1.55cm]current page.north east)
    {\sffamily\bfseries\fontsize{9}{12}\selectfont I.E. Sor María Juliana\\Cartago, Valle del Cauca};

  \node[anchor=west,text=milcGradeText] (gradonum) at ([xshift=1.85cm,yshift=-7.1cm]current page.north west)
    {\sffamily\bfseries\fontsize{112}{112}\selectfont 8};
  % El circulito de grado (°) solo aplica a las guias de grado («11°»); en
  % Bebras y Semillero el numero grande es un momento/modulo, no un grado.
  % Se posiciona relativo al nodo `gradonum`, no en coordenadas absolutas.
  \draw[milcGradeText,line width=1.6pt]
    ([xshift=2.0mm,yshift=-4.5mm]gradonum.north east) circle (.15cm);

  \node[anchor=west,text=milcGradeText,text width=11.4cm,align=left] at ([xshift=1.5cm,yshift=0.5cm]gradonum.east)
    {
     {\sffamily\bfseries\fontsize{9}{11}\selectfont\MakeUppercase{Período 2 · Lógica, algoritmos y computación física}}\\[3mm]
     {\sffamily\bfseries\fontsize{21}{25}\selectfont\setlength{\baselineskip}{27pt}Buscar el error\\[4pt]de a una puntada}\\[3.5mm]
     {\sffamily\bfseries\fontsize{15}{18}\selectfont Guía 17-8-TIC}
    };
\end{tikzpicture}

\vspace*{9.4cm}
\vfill

{\sffamily\bfseries\fontsize{8}{10}\selectfont\textcolor{milcGradeDark}{LO QUE TE LLEVAS HOY}}

\begin{tcolorbox}[enhanced,colback=white,colframe=milcNegro,arc=2mm,boxrule=1.6pt,
  left=5mm,right=5mm,top=4mm,bottom=4mm,before skip=3pt,after skip=12pt,
  fontupper=\RaggedRight\fontsize{11}{15.5}\selectfont]
Tres errores corregidos en un programa de MakeCode que te entrega tu docente o tu pareja, con una bitácora de depuración de tres columnas por cada uno: síntoma, hipótesis y prueba con su resultado. Y una nota final de cinco líneas: cuál error fue el más difícil, por qué, y qué técnica usaste para encontrarlo.
\end{tcolorbox}

\begin{tcolorbox}[enhanced,colback=milcGris,colframe=milcGris,arc=2mm,boxrule=0pt,
  left=5mm,right=5mm,top=3.5mm,bottom=3.5mm,after skip=12pt]
{\sffamily\bfseries\fontsize{8}{10}\selectfont\textcolor{milcNegro!55}{ESTA GUÍA ES DE}}\\[3mm]
\begin{tabularx}{\linewidth}{X p{3.2cm} p{3.2cm}}
\linead{\linewidth} & \linead{3.0cm} & \linead{3.0cm}\\[-1mm]
{\fontsize{8.5}{10}\selectfont\textcolor{milcNegro!55}{Nombre completo}} &
{\fontsize{8.5}{10}\selectfont\textcolor{milcNegro!55}{Grupo}} &
{\fontsize{8.5}{10}\selectfont\textcolor{milcNegro!55}{Fecha}}\\
\end{tabularx}
\end{tcolorbox}

\vfill

\noindent\textcolor{milcLinea}{\rule{\linewidth}{1pt}}\\[2mm]
{\sffamily\bfseries\fontsize{9.5}{12}\selectfont Autor: Dr. Álvaro Cárdenas Orozco}\\[1mm]
{\sffamily\fontsize{8.5}{11}\selectfont\textcolor{milcNegro!55}{Metodología MILC · apertura ancestral · pensamiento computacional · triángulo Dussel–estoicismo–Floridi}}

\newpage

\section*{Apertura: Depuración --- una hipótesis, una prueba}

\begin{softbox}{milcGradeDark}{milcGradeSoft}{Saber ancestral}
En Cartago, una prenda bordada no la hace una sola persona. Pasa por unas ocho artesanas y más de ocho horas de trabajo antes de llevar el nombre de un taller (Chica García, 2019). Cada una recibe la pieza de la anterior y la revisa antes de seguir: si una puntada quedó floja o un hilo cambió de tono, se ve en esa mesa y se corrige antes de que la prenda avance. Nadie deshace todo el bordado para encontrar una puntada mal hecha. Se busca por partes: esta flor, este tallo, esta franja. Los ornamentos que vistió el papa en su visita a Colombia salieron de esas mesas. Un programa con errores se revisa igual. No se cambia todo a la vez: se mira una parte, se prueba, se descarta o se confirma, y se pasa a la siguiente. La \textbf{cara de exclusión}: el nombre del taller lo sostienen ocho mujeres que rara vez aparecen en él; la reputación se acumula en el taller, no en quien borda. Hoy vas a buscar errores como se revisa una prenda: de a una parte, con una hipótesis y una prueba cada vez.\par\smallskip{\footnotesize\textcolor{milcNegro!70}{\textbf{Fuente.} Chica García, A. (2019, 17 de agosto). \emph{Bordados de Cartago: la herencia española que apropiaron las mujeres vallunas}. Radio Nacional de Colombia.}}
\end{softbox}

\begin{softbox}{milcTurquesa}{milcGris}{Saber contemporáneo}
Un \textbf{error} de programa, en inglés \textit{bug}, es una pieza que no hace lo que esperabas: una condición al revés, una variable que nunca se llenó, una pausa que falta. \textbf{Depurar} es encontrarlo y corregirlo. Quien programa pasa más tiempo buscando errores que escribiendo bloques nuevos, y por eso el método importa. Tiene cuatro pasos. (1) \textbf{Reproducir}: saber qué entrada provoca el error, siempre, no a veces. (2) \textbf{Hipótesis}: dos o tres explicaciones posibles, la más probable primero. (3) \textbf{Una prueba}: cambias una sola cosa y miras qué pasa. (4) \textbf{Confirmar o descartar}: si se arregló, era eso; si no, vuelves al programa original y pruebas la siguiente. Tres técnicas ayudan: mostrar el valor de una variable en pantalla, deshabilitar un bloque para aislar la parte que falla, y probar con valores extremos.
\end{softbox}

\begin{softbox}{milcMostaza}{milcGris}{Pregunta-puente}
\textit{La bordadora revisa esta flor, no toda la prenda. Cuando tu micro:bit muestre una cara feliz con 10 grados, ¿por dónde empiezas a mirar: por todo el programa o por el bloque que compara la temperatura?}
\end{softbox}

\begin{softbox}{milcVerde}{milcGris}{Saber y saber hacer}
Al terminar podrás: (1) \textbf{identificar} tres hipótesis sobre un error sin tocar el código; (2) \textbf{aplicar} el método de una prueba a la vez con bitácora; (3) \textbf{evaluar} cuál técnica de depuración sirvió con cada error; (4) \textbf{explicar} por qué cambiar varias cosas a la vez no es depurar.
\end{softbox}



\small
\begin{tabularx}{\textwidth}{>{\bfseries}p{3.0cm}Y}
\rowcolor{milcGradePrimary!12}Autor & Dr. Álvaro Cárdenas Orozco\\
DBA & Pensamiento computacional aplicado a sistemas de control con sensores y actuadores (MEN, T\&I)\\
Referentes & Pensamiento computacional · MakeCode / micro:bit · Logos como razón universal\\
Producto final & Tres errores corregidos en un programa de MakeCode que te entrega tu docente o tu pareja, con una bitácora de depuración de tres columnas por cada uno: síntoma, hipótesis y prueba con su resultado. Y una nota final de cinco líneas: cuál error fue el más difícil, por qué, y qué técnica usaste para encontrarlo.\\
\end{tabularx}
\normalsize

\puente{En la sesión 6 calibraste umbrales y viste que un programa puede fallar por un número mal puesto. Hoy aprendes a encontrar el error con método, sin adivinar. En la sesión 8 vas a combinar sensores con lógica compuesta, y la depuración de hoy será tu red.}

\milcestacion{milcGradeDark}{Ruta MILC: así vamos a aprender}
\begin{tabularx}{\textwidth}{>{\centering\arraybackslash}X>{\centering\arraybackslash}X>{\centering\arraybackslash}X>{\centering\arraybackslash}X}
\phasecard{1}{milcVerde}{milcGris}{Escucha\\[-1mm]\small Observo, escucho y pregunto.} &
\phasecard{2}{milcTurquesa}{milcGris}{Entender\\[-1mm]\small Sistematizo y ordeno.} &
\phasecard{3}{milcMagenta}{milcGris}{Hacer\\[-1mm]\small Construyo y aplico.} &
\phasecard{4}{milcVino}{milcGris}{Cierre\\[-1mm]\small Evalúo y reflexiono.}
\end{tabularx}


\puente{Antes de corregir nada, vas a mirar un programa con un error y escribir tres hipótesis. Es la mirada de la bordadora sobre la prenda: primero ver, después tocar.}

\milcestacion{milcVerde}{Estación 1 de 3 · Escucha}

\begin{softbox}{milcVerde}{milcGris}{Escena para escuchar}
\textbf{Actividad 1 · IDENTIFICA --- Tres hipótesis sin tocar el código} (15 min · individual). (1) Tu docente proyecta un programa de MakeCode con un error: debería mostrar una cara triste cuando la temperatura baja de 18 grados, pero muestra la cara feliz. (2) Sin tocar nada, mira el bloque «si… si no» y copia en el cuaderno la condición tal como está. (3) Escribe tres hipótesis: ¿la comparación está al revés?, ¿los íconos están en la rama equivocada?, ¿la variable de temperatura nunca se llena? (4) Ordénalas de la más probable a la menos probable y escribe por qué pusiste esa primero. (5) Escribe qué prueba harías para la primera hipótesis, y solo para esa.
\end{softbox}

{\sffamily\bfseries\fontsize{8}{10}\selectfont\textcolor{milcNegro!55}{MARCA CUANDO LO HAYAS HECHO}}
\milccheckrow{Copié la condición del programa tal como está.}
\milccheckrow{Escribí tres hipótesis y las ordené con una razón.}
\milccheckrow{Escribí una sola prueba para la primera hipótesis.}

\begin{infoband}{milcVerde}


\textbf{Lo que va al cuaderno (Actividad 1 · IDENTIFICA).} \textbf{Título:} «Actividad 1 --- Tres hipótesis sin tocar el código». \textbf{Formato:} la condición copiada, las tres hipótesis numeradas con la razón del orden, y la prueba para la primera. \textbf{Extensión:} media página. \textbf{Sabes que terminaste cuando} tu prueba cambia una sola cosa y puedes decir qué esperas ver si la hipótesis es cierta.
\end{infoband}

\puente{Ya tienes tres hipótesis ordenadas. Ahora aprendes las tres técnicas y, con tu pareja, corriges el primer error con una prueba a la vez.}

\milcestacion{milcTurquesa}{Estación 2 de 3 · Entender}

\begin{softbox}{milcTurquesa}{milcGris}{Lo que vas a entender}
\textbf{Actividad 2 · APLICA --- El primer error, con bitácora} (30 min · parejas). (1) Tu docente entrega un programa con tres errores, o intercambias con tu pareja un programa donde cada uno sembró tres. (2) Con tu pareja, pruébenlo en el simulador hasta que el primer error aparezca siempre con la misma entrada: eso es reproducirlo. (3) Dibujen la bitácora de tres columnas: síntoma, hipótesis, prueba y resultado. (4) Escriban el síntoma y dos hipótesis; prueben solo la primera cambiando una sola cosa. (5) Si se arregló, anoten «confirmada»; si no, vuelvan al programa original, anoten «descartada» y prueben la segunda.
\end{softbox}

\milcpaso{1}{milcTurquesa}{\textbf{Reproducir antes que nada.} Un error que aparece «a veces» no se puede buscar. Prueba con la misma entrada, 10 grados, botón A, hasta que falle siempre. Cuando sabes qué lo provoca, ya tienes la mitad.}
\milcpaso{2}{milcTurquesa}{\textbf{Las tres técnicas.} (1) Mostrar la variable: pon «mostrar número (temperatura)» antes del «si» y mira qué valor llega. (2) Aislar: clic derecho sobre un bloque y «deshabilitar»; si el error desaparece, estaba ahí. (3) Valores extremos: prueba con 0 y con 40 grados en el simulador; los extremos muestran lo que el valor normal esconde.}
\milcpaso{3}{milcTurquesa}{\textbf{Una cosa a la vez.} Si cambias la comparación y el ícono en la misma prueba y se arregla, no sabes cuál de los dos era. Y si se daña más, tampoco. Cada prueba cambia una sola cosa, y la bitácora dice cuál.}
\milcpaso{4}{milcTurquesa}{\textbf{Cómo se hace, paso a paso.} (1) Reproducir con la misma entrada. (2) Escribir el síntoma. (3) Dos hipótesis, la más probable primero. (4) Una prueba, un cambio. (5) Confirmar o volver atrás y probar la siguiente. (6) Anotar todo.}

\begin{drawbox}{milcTurquesa}{La depuración, en una mirada}
\textbf{Síntoma:} qué hace mal y con qué entrada. \\[1mm] \textbf{Hipótesis:} dónde crees que está el error. \\[1mm] \textbf{Prueba:} el único cambio que haces. \\[1mm] \textbf{Resultado:} confirmada o descartada. \\[1mm] \textbf{Técnica:} mostrar variable, aislar o valores extremos.
\end{drawbox}

\begin{softbox}{milcMostaza}{milcGris}{Errores comunes y su corrección}
(1) \textbf{Cambiar varias cosas a la vez}: no sabes cuál arregló ni cuál dañó. (2) \textbf{No reproducir}: buscas un error que no sabes cuándo aparece. (3) \textbf{Suponer sin probar}: «debe ser la variable» y cambias sin mirar. (4) \textbf{Sin bitácora}: repites la misma prueba tres veces sin darte cuenta. (5) \textbf{Mover bloques al azar}: cuando el error no cede, el programa termina peor.
\end{softbox}

\begin{infoband}{milcTurquesa}


\textbf{Lo que va al cuaderno (Actividad 2 · APLICA).} \textbf{Título:} «Actividad 2 --- El primer error, con bitácora». \textbf{Formato:} la bitácora de tres columnas con la entrada que reproduce el error, el síntoma, las hipótesis y cada prueba con su resultado. \textbf{Extensión:} media página. \textbf{Sabes que terminaste cuando} el primer error está corregido y la bitácora dice qué cambiaste en cada prueba y qué pasó.
\end{infoband}

\puente{Ya corregiste un error con bitácora. Quedan dos, y en cada uno vas a anotar qué técnica te sirvió.}

\milcestacion{milcMagenta}{Estación 3 de 3 · Hacer}

\begin{softbox}{milcMagenta}{milcGris}{Cómo vas a construir tu producto}
\textbf{Actividad 3 · EVALÚA --- Dos errores más y la técnica que sirvió} (25 min · individual). (1) Reproduce el segundo error y llena su bitácora: síntoma, hipótesis, una prueba, resultado. (2) Haz lo mismo con el tercero. (3) Junto a cada error corregido, escribe qué técnica te sirvió: mostrar variable, aislar o valores extremos. (4) Escribe la nota final de cinco líneas: cuál error fue el más difícil y por qué. (5) Pásale tu bitácora a tu pareja y pídele que diga, sin verte el programa, en qué orden probaste.
\end{softbox}

\milcpaso{1}{milcMagenta}{\textbf{Tu entrega tiene tres piezas.} (1) El programa con los tres errores corregidos, compartido con enlace o archivo .hex. (2) La bitácora de tres columnas por cada error. (3) La nota final sobre el más difícil y la técnica que usaste.}
\milcpaso{2}{milcMagenta}{\textbf{Los tres errores típicos.} La comparación al revés: «mayor que» donde iba «menor que». La variable vacía: se usa antes de haberla llenado con «establecer». La pausa que falta: la animación corre tan rápido que solo se ve el último cuadro. Cada uno se ve con una técnica distinta.}
\milcpaso{3}{milcMagenta}{\textbf{Descartar también es avanzar.} Una hipótesis descartada con una prueba limpia te deja con una menos. Anótala igual que la confirmada. La bordadora que revisó la flor y la encontró bien no perdió el tiempo: ya sabe que el problema está en otra parte.}
\milcpaso{4}{milcMagenta}{\textbf{La prueba del que lee.} Tu bitácora sirve si otra persona sabe, sin ver el programa, qué probaste y en qué orden. Si tu pareja no puede reconstruirlo, faltó anotar el cambio o el resultado.}

\begin{drawbox}{milcMagenta}{Antes de entregar}
\checkbox Cada error tiene su entrada que lo reproduce. \\[1mm] \checkbox Cada error tiene al menos una hipótesis escrita antes de probar. \\[1mm] \checkbox Cada prueba cambia una sola cosa. \\[1mm] \checkbox Bitácora de tres columnas por cada uno de los tres errores. \\[1mm] \checkbox Los tres errores corregidos y probados. \\[1mm] \checkbox Nota final con el error más difícil y la técnica que sirvió. \\[1mm] \checkbox Programa compartido con enlace o archivo .hex.
\end{drawbox}

\begin{softbox}{milcMagenta}{milcGris}{Plantilla de trabajo}
\textbf{Síntoma.} \textit{Con … grados muestra …} \\[1mm] \textbf{Hipótesis 1.} \textit{Creo que …} \\[1mm] \textbf{Prueba.} \textit{Cambié solo …} \\[1mm] \textbf{Resultado.} \textit{Confirmada / descartada porque …} \\[1mm] \textbf{Técnica.} \textit{Mostrar variable / aislar / valores extremos.}
\end{softbox}

\begin{infoband}{milcMagenta}


\textbf{Lo que va al cuaderno (Actividad 3 · EVALÚA).} \textbf{Título:} «Actividad 3 --- Dos errores más y la técnica que sirvió». \textbf{Formato:} las bitácoras del segundo y el tercer error, la técnica junto a cada uno, y la nota final de cinco líneas. \textbf{Extensión:} una página. \textbf{Sabes que terminaste cuando} los tres errores están corregidos y tu pareja pudo decir en qué orden probaste solo con leer tu bitácora.
\end{infoband}

\puente{Lo que entregas hoy: tres errores corregidos, la bitácora de tres columnas por cada uno y la nota sobre el más difícil. La prueba de calidad: alguien lee tu bitácora y sabe qué probaste, en qué orden y qué pasó.}

\milcestacion{milcMostaza}{Producto de la sesión}
\textbf{Tres errores corregidos + bitácora de tres columnas por error + nota sobre el más difícil y la técnica usada}.
\begin{softbox}{milcMostaza}{milcGris}{Criterios de calidad}
(1) Cada error tiene la entrada que lo reproduce y su síntoma escrito. (2) Cada prueba de la bitácora cambia una sola cosa y tiene resultado. (3) Los tres errores están corregidos y el programa funciona con las entradas de prueba. (4) Junto a cada error está la técnica que sirvió. (5) La nota final dice cuál fue el más difícil y por qué. (6) El programa está compartido con enlace o archivo .hex.
\end{softbox}

\puente{Antes de cerrar, mira lo que hiciste desde cinco lados. Un error encontrado con método se puede contar; uno que se arregló solo no enseña nada.}

\milcestacion{milcVino}{Evaluación liberadora · cinco dimensiones}

\begin{softbox}{milcVino}{milcGris}{Sentido de la evaluación}
Evaluar no es solo revisar una nota. Es reconocer qué aprendí, qué cambió en mí, qué emociones manejé, cómo me relaciono con la ciudad y la comunidad, y qué le dejaré a quien viene detrás.
\end{softbox}

\textbf{Mi reflexión liberadora en cinco dimensiones.} Completa cada frase iniciada:

\small
\begin{tabularx}{\textwidth}{>{\bfseries\RaggedRight\arraybackslash}p{3.4cm}Y>{\RaggedRight\arraybackslash}p{4.6cm}}
\rowcolor{milcVino}\color{white}Dimensión & \color{white}\bfseries Mi respuesta (frame starter) & \color{white}\bfseries Algo que descubrí\\
1. Personal & Lo que más cambió en mí fue \linead{6.5cm} & \linead{4.2cm}\\
2. Emocional & La emoción más fuerte fue \linead{2.5cm} y la regulé \linead{3cm} & \linead{4.2cm}\\
3. Ciudadana & En democracia esto importa porque \linead{6cm} & \linead{4.2cm}\\
4. Local (Cartago) & En mi barrio sería útil para \linead{6.5cm} & \linead{4.2cm}\\
5. Intergeneracional & A mi abuela/abuelo le diría \linead{6.5cm} & \linead{4.2cm}\\
\end{tabularx}
\normalsize

\begin{infoband}{milcVino}
\textbf{Para la próxima sustentación.} Vuelve a esta tabla en seis meses. Verás que las respuestas cambian — no porque lo aprendido cambie, sino porque tú cambias. La evaluación liberadora es una conversación contigo a través del tiempo.
\end{infoband}


\puente{Para terminar, tres ideas sobre mirar antes de suponer, contadas con palabras de esta guía: Dussel sobre el silencio del que escucha, Marco Aurelio sobre cambiar de parecer, y la Onlife Initiative sobre lo que rechazamos porque no entendemos.}

\milcestacion{milcGradeDark}{Tres ideas para cerrar · Dussel, un estoico y Floridi}

\begin{softbox}{milcVino}{milcGris}{Enrique Dussel · lente del nosotros}
\textit{El respeto es el silencio de quien todo tiene que escuchar, porque todavía no sabe nada del otro.}

{\footnotesize\itshape\color{milcNegro!70}--- idea tomada de Enrique Dussel, contada con palabras de esta guía. No es una cita textual.}

Antes de cambiar un bloque, el programa te está diciendo algo con su síntoma. Reproducir y escribir hipótesis es escuchar primero. Mover bloques al azar es hablar sin haber oído.

\textbf{¿Qué problema de esta semana intenté arreglar antes de entender qué pasaba?}
\end{softbox}
\linead{\linewidth}\\[1mm]
\linead{\linewidth}

\begin{softbox}{milcMostaza}{milcGris}{Estoicismo · Marco Aurelio · cuidado interior}
\textit{Cambiar de parecer cuando alguien te corrige no te quita libertad: es parte de ella.}

{\footnotesize\itshape\color{milcNegro!70}--- idea tomada de Marco Aurelio, contada con palabras de esta guía. No es una cita textual.}

Tu hipótesis favorita puede estar equivocada. La prueba que la descarta te está corrigiendo. Anotar «descartada» y pasar a la siguiente es lo que hace quien depura; insistir es lo que hace quien adivina.

\textbf{¿Cuál hipótesis me costó soltar hoy aunque la prueba la había descartado?}
\end{softbox}
\linead{\linewidth}\\[1mm]
\linead{\linewidth}

\begin{softbox}{milcTurquesa}{milcGris}{Luciano Floridi y la Onlife Initiative · lente de la infoesfera}
\textit{Tememos y rechazamos aquello a lo que no logramos darle sentido.}

{\footnotesize\itshape\color{milcNegro!70}--- idea tomada de Luciano Floridi y la Onlife Initiative, contada con palabras de esta guía. No es una cita textual.}

Un error que aparece «a veces» da miedo porque no tiene sentido todavía. Reproducirlo con la misma entrada le da sentido, y un error con sentido ya no asusta: se corrige.

\textbf{¿Qué cosa de la tecnología rechazo porque no la entiendo, y qué prueba me ayudaría a entenderla?}
\end{softbox}
\linead{\linewidth}\\[1mm]
\linead{\linewidth}

\begin{infoband}{milcNegro}
\textbf{Nota para el docente.} Las tres ideas están contadas con palabras de esta guía a partir del pensamiento de cada autor; no son citas textuales. De 9.º en adelante se trabajan con la obra y su referencia.
\end{infoband}

\milcestacion{milcVino}{Mi autoevaluación · marca una casilla por fila}

{\small
\setlength{\extrarowheight}{2.2pt}
\begin{tabularx}{\textwidth}{>{\RaggedRight\arraybackslash\bfseries}p{3.0cm}YYY}
\rowcolor{milcVino}\color{white}\bfseries Criterio & \color{white}\bfseries Lo logré & \color{white}\bfseries En proceso & \color{white}\bfseries Necesito apoyo\\[.6mm]
Comprensión del tema & \milcopcion{Explico las ideas con mis palabras.} & \milcopcion{Reconozco las ideas, falta precisión.} & \milcopcion{Necesito repasar lo básico.}\\[.8mm]
\rowcolor{milcGris}Pensamiento computacional & \milcopcion{Usé los pilares al armar mi producto.} & \milcopcion{Usé algunos pilares.} & \milcopcion{Necesito releer la estación 2.}\\[.8mm]
Aplicación y producto & \milcopcion{Mi producto cumple los criterios.} & \milcopcion{Está armado, con ajustes pendientes.} & \milcopcion{Aún está en construcción.}\\[.8mm]
\rowcolor{milcGris}Reflexión 5 dimensiones & \milcopcion{Respondí cada una con honestidad.} & \milcopcion{Respondí algunas con detalle.} & \milcopcion{Mis respuestas son muy generales.}\\[.8mm]
Triángulo de pensamiento & \milcopcion{Conecté las 3 voces con mi trabajo.} & \milcopcion{Conecté al menos una.} & \milcopcion{No las relaciono con mi proceso.}\\[.8mm]
\end{tabularx}}

\vspace{3mm}
\textbf{Hoy entendí que:}\\[1mm]
\linead{\linewidth}\\[1mm]
\linead{\linewidth}

\textbf{Mi compromiso para la próxima sesión es:}\\[1mm]
\linead{\linewidth}\\[1mm]
\linead{\linewidth}

\begin{softbox}{milcMostaza}{milcGris}{Cita de despedida · Marco Aurelio}
\textit{``El obstáculo en el camino se vuelve el camino. Lo que estorba al camino se vuelve camino.''}\\[1mm]
Cada error de hoy, bien observado, se convierte en pieza del camino que pisarás mañana.
\end{softbox}

\small
\begin{tabularx}{\textwidth}{>{\bfseries}p{3.6cm}Y}
\rowcolor{milcGradeDark!12}Nombre completo: & \linead{10cm}\\
Fecha: & \linead{4cm} \quad Curso/grupo: \linead{4cm}\\
Mi firma: & \linead{9cm}\\
\end{tabularx}
\normalsize

\begin{infoband}{milcGradeDark}
\textbf{Para la próxima sesión.} Llega con tu bitácora de los tres errores. En la sesión 8 vas a construir alertas que combinan sensores con Y, O y NO. Las vas a probar con una tabla de ocho escenarios. Cuando una fila no coincida, vas a usar el método de hoy.
\end{infoband}

\end{document}
